\section{Born 权重的等价表述：由计数概率构造有效 Hilbert 描述}

\subsection{从分布到幅度：一个定义级构造}

\begin{definition}[有效 Hilbert 空间与基]
\label{def:effective_hilbert_space}
固定 $m$。令
$$
\mathcal H_m:=\CC^{X_m}
$$
并赋予标准内积 $\langle f,g\rangle=\sum_{x\in X_m}\overline{f(x)}g(x)$。对每个 $x\in X_m$ 定义基向量 $|x\rangle$ 为指示函数：
$|x\rangle(y)=\mathbf 1_{y=x}$。
\end{definition}

\begin{definition}[由计数概率诱导的态向量族]
\label{def:psi_from_p}
给定计数概率 $p_m:X_m\to[0,1]$，选取任意相位函数 $\theta:X_m\to\RR$，定义
$$
|\psi_m\rangle:=\sum_{x\in X_m}\sqrt{p_m(x)}\,e^{i\theta(x)}|x\rangle\in\mathcal H_m.
$$
\end{definition}

\begin{proposition}[归一性]
\label{prop:psi_normalized}
若 $\sum_{x\in X_m}p_m(x)=1$，则 $\langle\psi_m|\psi_m\rangle=1$。
\end{proposition}

\subsection{路线 A 的核心结论：Born 权重仅是计数概率的另一种写法}

\begin{theorem}[Born 权重的计数等价表述]
\label{thm:born_as_counting_equivalence}
对任意 $x\in X_m$，令 $P_x:=|x\rangle\langle x|$ 为基投影算子，则
$$
\Pr(x)
:=\langle \psi_m|P_x|\psi_m\rangle
=|\langle x|\psi_m\rangle|^2
=p_m(x).
$$
\end{theorem}

\begin{proof}
由定义 $\langle x|\psi_m\rangle=\sqrt{p_m(x)}e^{i\theta(x)}$，故其模平方等于 $p_m(x)$。又 $\langle\psi_m|P_x|\psi_m\rangle=\langle\psi_m|x\rangle\langle x|\psi_m\rangle=|\langle x|\psi_m\rangle|^2$。
\end{proof}

\begin{remark}[相位的地位与开放项]
定理 \ref{thm:born_as_counting_equivalence} 仅涉及基投影读出，因此相位 $\theta(x)$ 不影响 $\Pr(x)$。
相位信息只有在引入更一般的效应（非对角元）或更细的读出语义时才变得可见。
在折叠本体中，一个自然的候选是用离散连接/holonomy 结构为相位提供几何落点；但这需要额外输入（例如具体的平行移动实现、可干涉读出协议等），本文将其列为讨论节的开放方向。
\end{remark}

